\documentclass[11pt]{extarticle}

\usepackage[warn]{mathtext}
\usepackage[T2A]{fontenc}      
\usepackage[utf8]{inputenc}      
\usepackage[english,russian]{babel} 
\usepackage{indentfirst}
\usepackage{csquotes}
% \usepackage[bibstyle=gost-numeric, sorting=none]{biblatex}
% \addbibresource{biblio.bib}

\usepackage[
    left=1.8cm,
    right=1.8cm,
    top=1.8cm,
    bottom=1.8cm,
    bindingoffset=0cm
]{geometry}

\usepackage{graphicx, hyperref, xcolor}
\hypersetup{
    colorlinks=true,
    linkcolor=teal,
    filecolor=magenta, 
    urlcolor=blue,
    citecolor=olive,
    pdftitle={GD},
    % pdfpagemode=FullScreen,
    linktoc=all
    }

\usepackage{wrapfig,caption}

% figures
\usepackage{caption}
\usepackage{subcaption}
\usepackage{floatrow}
\floatsetup{heightadjust=object}

% math
\usepackage{amsmath,amsfonts,amssymb,amsthm,mathtools,esint,eucal}

\begin{document}
\section*{Аннотация к статье: \\ "Reward Tampering Problems and Solutions in Reinforcement Learning: \\ A Causal Influence Diagram Perspective"}

Авторы статьи: Tom Everitt, Marcus Hutter, Ramana Kumar, Victoria Krakovna

Организация: DeepMind, Australian National University

Дата выпуска: March 29, 2021

\section*{Ключевые идеи}
Статья посвящена фундаментальной проблеме обучения с подкреплением — \emph{вмешательству в функцию вознаграждения} (reward tampering). Авторы показывают, что агенты, стремящиеся максимизировать наблюдаемое вознаграждение, могут не решать поставленную задачу, а манипулировать системой наград. В статье предлагается формальная классификация видов вмешательства: \emph{вмешательство в функцию вознаграждения} (Reward Function Tampering), \emph{вмешательство во входы функции вознаграждения} (RF-Input Tampering), и \emph{вмешательство в процесс обучения вознаграждения} (Reward Learning Tampering).

Главная цель работы — разработать принципы проектирования агентов, при которых вмешательство в функцию вознаграждения перестаёт быть оптимальной стратегией. Среди этих принципов: \emph{counterfactual rewards}, \emph{uninfluenceable learning}, \emph{current-RF optimization}, \emph{history-based rewards} и \emph{belief-based rewards}.


\section*{Mатематическое описание метода}

Среда формализуется как частично наблюдаемая марковская процесс (POMDP) с состоянием $S_t$, действием $A_t$, наблюдением $O_t$ и вознаграждением $R_t$.

Основная динамика описывается как

\begin{equation}
P(S_{t+1}, O_t, R_t \mid S_t, A_t, \Theta^R, \Theta^T, \Theta^O)
\end{equation}

где $\Theta^R$ ~-- параметры функции вознаграждения, $\Theta^T$ ~-- параметры функции переходов среды, а $\Theta^O$ ~-- параметры функции наблюдения.

Задача агента формулируется как максимизация математического ожидания суммарного вознаграждения:

\begin{equation}
V^\pi = \mathbb{E}\left[\sum_{t=1}^T R_t \mid \pi, \Theta^R, \Theta^T, \Theta^O \right]
\end{equation}

Проблема вмешательства возникает, когда существует причинный путь
$A_t \to \Theta^R \to R_{t+1}$,
т.е. действия агента могут изменять функцию вознаграждения. Чтобы устранить эту зависимость, предлагается использовать \emph{временную изоляцию} (temporal isolation, TI) и \emph{неподверженность влиянию} (uninfluenceability), которые блокируют этот путь в причинной структуре.

Для случая частичного наблюдения используется состояние убеждений $B_t$, определяемое как аппроксимация апостериорного распределения $P(S_t \mid O_{1:t}, A_{1:t-1})$. Тогда вознаграждение может быть задано как функция от этого состояния:

\begin{equation}
R_t = R(B_t; \Theta^R)
\end{equation}

что предотвращает прямое вмешательство агента в наблюдения или функцию вознаграждения.

\section*{Алгоритмы}

\paragraph{Базовый RL-агент (референс, Прил. C).} Агент с фиксированной, неизменяемой функцией вознаграждения оптимизирует ожидаемую сумму наград по известной $\Theta^R$:
\begin{equation}
\pi^*=\arg\max_{\pi},\mathbb{E}\!\left[\sum_{t=1}^{m} R(S_t;\Theta^R)\right].
\end{equation}

\paragraph{Current-RF Optimization (CRFO), TI-considering.} Агент на шаге $k$ планирует по текущей реализованной RF и учитывает возможные будущие изменения $\hat{\Theta}^R$; у него может возникать цель сохранять текущую RF. Учитывание TI делает агенту выгодным избегать изменений $\hat{\Theta}^R$ при выполнении предпосылок приватности или неинформативности. Целевая функция на шаге $k$:
\begin{equation}
\pi_k=\arg\max_{\pi},\mathbb{E}\!\left[\sum_{t=k}^{m} R(S_t;\hat{\Theta}^R_k)\right].
\end{equation}


\paragraph{Current-RF Optimization (CRFO), TI-ignoring.} Агент на шаге $k$ планирует как будто будущие действия тоже будут выбраны относительно текущей $\hat{\Theta}^R_k$, т.е. игнорирует будущие сдвиги RF; это убирает стимул влиять на будущую RF:
\begin{equation}
A_k=\arg\max_{a_k},\mathbb{E}\!\left[\sum_{t=k}^{m} R(S_t;\hat{\Theta}^R_k)\mid S_k,;A_k=a_k\right].
\end{equation}

\paragraph{Uninfluenceable Learning — Direct Learning.} Агент байесовски выводит целевую RF пользователя по данным обратной связи $D_{1:t}$. При приватности или неинформативности данных единственная инструментальная цель — сделать $D_t$ более информативными о $\Theta^R_*$. Процесс обучения сконструирован так, чтобы его результат не зависел от действий агента:
\begin{equation}
P(\Theta^R_*\mid D_{1:t})\propto P(D_{1:t}\mid \Theta^R_*),P(\Theta^R_*),\qquad
\pi^*=\arg\max_{\pi},\mathbb{E}\!\left[\sum_{t=1}^{m} R(S_t;\Theta^R_*)\right].
\end{equation}

\paragraph{Uninfluenceable Learning — Counterfactual RF.} Агент оптимизирует \emph{контрфактическую} реализованную RF, которую получил бы процесс обучения при фиксированной безопасной политике. Таким образом, инструментальная цель ограничивается повышением информативности об \emph{контрфактической} RF, а не воздействием на реальную:

$\pi_{\mathrm{safe}}$ (twin-network):
\begin{equation}
\tilde{\Theta}^R_t = \mathcal{F}!\left(D^{\pi_{\mathrm{safe}}}{1:t}\right),\qquad
\pi^*=\arg\max{\pi},\mathbb{E}\!\left[\sum_{t=1}^{m} R(S_t;\tilde{\Theta}^R_t)\right].
\end{equation}


\paragraph{RF-Input Tampering — History-Based Rewards.} Вознаграждение определяется по полной истории наблюдений и действий, что позволяет сделай корректную политику или критерий достижения цели и исключить выгоду от подделки текущих наблюдений:
\begin{equation}
R_t = R\big(O_{1:t},A_{1:t}\big),\qquad
\pi^*=\arg\max_{\pi},\mathbb{E}\!\left[\sum_{t=1}^{m} R\big(O_{1:t},A_{1:t}\big)\right].
\end{equation}

\paragraph{RF-Input Tampering — Belief-Based Rewards.} Вознаграждение задаётся на состоянии убеждений $B_t$ предсказательной модели $\Theta^{PM}$; при TI-ignoring это убирает стимулы править будущую модель:
\begin{equation}
R_t = R(B_t;\Theta^R),\qquad
V^\pi=\mathbb{E}\!\left[\sum_{t=1}^{m} R(B_t;\Theta^R)\mid B_1,\Theta^{PM}\right],\quad \pi^*=\arg\max_{\pi} V^\pi.
\end{equation}


\section*{Наборы данных}
Работа является теоретической и аналитической и не использует эмпирические наборы данных. Для иллюстрации вводятся гипотетические примеры сред, такие как «камни и алмазы» (rocks and diamonds), где агент может собирать объекты, либо искажать свои наблюдения для повышения награды. Эти примеры демонстрируют, как различные подходы предотвращают вмешательство во входы функции вознаграждения.
\section*{Метрики}
Поскольку цель статьи ~-- формальный анализ причинных структур, метрики производительности не используются.
Основные критерии — \emph{наличие или отсутствие инструментальной цели вмешательства}, то есть проверка того, существует ли причинный путь $A_t \to \Theta^R_{t+k} \to R_{t+k+1}$ в соответствующей диаграмме.
\section*{Модели}
Для формализации взаимодействия агент–среда используются \emph{диаграммы причинного влияния (Causal Influence Diagrams)}.
Каждый узел соответствует переменной ($S_t, A_t, O_t, R_t, B_t, \Theta^R, \Theta^{PM}$ и др.), а стрелки отражают причинные связи между ними.
\end{document}